% !TeX root = ../../Book1.tex
\section{切测度与局部放大}
\label{b1:sec:D101}

研究一张光滑曲面在一点附近的形状，可以先平移该点到原点，再把距离放大。
在合适的面积归一化下，曲面测度趋于切平面上的面积测度。
对于尚不知道是否可整流的测度，不能预先指定极限一定是平面；
应当先容许一般非零测度作为放大极限，再研究这些极限的结构。
这一区别使切测度既能描述已有的切平面，也能用于寻找切平面。

\subsection{归一化与非零极限}
\label{b1:subsec:D10101}

设 \(\mu\) 是 \(\mathbb R^n\) 上的正 Radon 测度，\(x\in\operatorname{supp}\mu\)。
对 \(r>0\)，记
\[
 \eta_{x,r}(y)=\frac{y-x}{r}.
\]
这里的推前不带 Jacobi 因子：
\[
 (\eta_{x,r})_\#\mu(A)=\mu(x+rA).
\]
尺度归一化另由一个正数给出，不能在推前的定义和前置系数中重复计算。

\begin{definition}\label{b1:def:app-d-tangent-measure}
若存在 \(r_j\downarrow0\) 和 \(c_j>0\)，使某个非零正 Radon 测度 \(\nu\) 满足
\[
 c_j(\eta_{x,r_j})_\#\mu\longrightarrow\nu
 \quad\text{在淡收敛意义下},
\]
则称 \(\nu\) 为 \(\mu\) 在 \(x\) 处的一个%
\term[qiecedu]{切测度}[tangent measure]。
所有这些测度组成的集合记为 \(\operatorname{Tan}(\mu,x)\)。
\symidx{\(\operatorname{Tan}(\mu,x)\)：测度在一点的切测度族}
\end{definition}

淡收敛由所有连续紧支集函数测试，沿用%
\cref{b1:def:vague-convergence}。
一般的切测度不是概率测度，定义中的 \(c_j\) 也不必等于某个预先指定的 \(r_j^{-k}\)。
排除零测度很重要：把 \(c_j\) 取得足够小，总能制造出没有几何信息的零极限。
另一方面，归一化不能自动保证紧性；放大后的质量还须在每个固定紧集上一致受控。

例如，若 \(\mu=\theta\mathcal H^k\llcorner P\)，其中 \(P\) 是经过 \(x\) 的 \(k\) 维仿射平面且 \(\theta>0\)，则
\[
 r^{-k}(\eta_{x,r})_\#\mu
   =\theta\mathcal H^k\llcorner(P-x).
\]
若 \(\mu\) 在 \(x\) 处有原子，固定系数放大产生的则是点质量，而非正维数的平面面积。
两个例子说明，尺度幂反映了测度在该点的维数，不能只根据环境空间的维数决定。

\begin{proposition}\label{b1:prop:app-d-tangent-scaling}
若 \(\nu\in\operatorname{Tan}(\mu,x)\)，则对任意 \(a,s>0\)，
\(a(\eta_{0,s})_\#\nu\) 也属于 \(\operatorname{Tan}(\mu,x)\)。
\end{proposition}
\begin{proof}
取定义中的 \(c_j,r_j\)。对 \(\varphi\in C_c(\mathbb R^n)\)，
\(\varphi\circ\eta_{0,s}\) 仍是连续紧支集函数，而且
\[
 \eta_{0,s}\circ\eta_{x,r_j}=\eta_{x,sr_j}.
\]
于是
\[
 ac_j(\eta_{x,sr_j})_\#\mu
   \longrightarrow a(\eta_{0,s})_\#\nu.
\]
极限非零，因为伸缩是同胚且 \(a>0\)。
\end{proof}

这条封闭性不表示每个切测度本身都齐次，也不表示切测度族对加法封闭。
“一个测度的所有伸缩仍在该族内”与“该测度在伸缩下只乘一个幂次”是不同断言。

\subsection{密度控制如何产生切测度}
\label{b1:subsec:D10102}

一般切测度理论需要细致选择尺度。本节先证明一个可直接使用的充分条件，
它已覆盖正有限密度点，也说明归一化和紧性各自承担什么工作。

\begin{proposition}\label{b1:prop:app-d-tangent-existence-density}
设 \(k>0\)。若存在 \(a,b,r_0>0\)，使
\[
 ar^k\leq\mu(B(x,r))\leq br^k
 \quad(0<r<r_0),
\]
则任意趋于零的尺度序列都有子列，使
\(r_j^{-k}(\eta_{x,r_j})_\#\mu\) 淡收敛于一个非零 Radon 测度。
\end{proposition}
\begin{proof}
记这些归一化测度为 \(\nu_j\)。
对每个固定 \(R>0\)，当 \(j\) 足够大时，
\[
 \nu_j(B(0,R))=r_j^{-k}\mu(B(x,Rr_j))\leq bR^k.
\]
有限个较早指标各自在这个球上也有有限质量。
因此 \((\nu_j)\) 在每个紧集上质量一致有界。
在依次增大的闭球上，对 \(C_c(\mathbb R^n)\) 的可数一致稠密测试族作对角抽取，
所得极限是局部有界的正线性泛函；由 Riesz 表示得到正 Radon 测度 \(\nu\)。
局部质量界把稠密测试族上的收敛推广到全部连续紧支集函数。

再取 \(0\leq\varphi\leq1\)，使 \(\varphi=1\) 于 \(\overline{B(0,1)}\)，
且 \(\operatorname{supp}\varphi\subset B(0,2)\)。
充分大的 \(j\) 满足 \(\int\varphi\,d\nu_j\geq a\)，
故 \(\int\varphi\,d\nu\geq a>0\)。
这一步使用下密度界，保证极限不因质量消失而退化为零。
\end{proof}

上密度界负责阻止固定观察窗口内质量无限增大，下密度界负责保留非零对象。
如果只写“把球的质量归一化为一”，仍须检查较大球上的质量；
一个单位球中的概率归一化并不控制所有固定半径的球。
本命题不声称是切测度存在的必要条件，也不替代一般 Radon 测度几乎处处的切测度理论。

\subsection{可整流测度的平面极限}
\label{b1:subsec:D10103}

\chapref{b1:ch:19}已给出可整流集和带权测度的密度结论。
把这些结论写成切测度语言，可以看清密度、权重和切平面如何同时进入极限。

\begin{theorem}\label{b1:thm:app-d-rectifiable-tangents}
设 \(1\leq k\leq n\)，\(E\subset\mathbb R^n\) 可数 \(k\) 维可整流，
\[
 \mu=\theta\mathcal H^k\llcorner E
\]
是 Radon 测度，其中 \(\theta\) 可测、非负。
则对 \(\mu\)-几乎每个 \(x\)，存在 \(k\) 维线性空间 \(P_x\)，使
\[
 r^{-k}(\eta_{x,r})_\#\mu
   \longrightarrow\theta(x)\mathcal H^k\llcorner P_x
 \quad(r\downarrow0),
\]
且
\[
 \operatorname{Tan}(\mu,x)
   =\{\,a\mathcal H^k\llcorner P_x:a>0\,\}.
\]
这里选取 \(\theta\) 的有限正值代表；结论只涉及满足该条件的满 \(\mu\) 测度集合。
\end{theorem}
\begin{proof}
先在 \(E\) 局部面积有限的情形工作。
由\cref{b1:thm:rectifiable-density-tangent}，在几乎每个 \(x\in E\)，
无权面积的归一化放大趋于 \(\mathcal H^k\llcorner P_x\)。
同时，由 Radon 测度的微分定理，几乎每个这样的点满足
\[
 \int_{E\cap B(x,r)}|\theta(y)-\theta(x)|\,d\mathcal H^k(y)
     =o(r^k).
\]
对支集包含于 \(B(0,R)\) 的测试函数 \(\varphi\)，有
\[
 \begin{split}
 &\left|r^{-k}\int_E
   \varphi\!\left(\frac{y-x}{r}\right)
   \bigl(\theta(y)-\theta(x)\bigr)\,d\mathcal H^k(y)\right|\\
 &\hspace{1em}\leq
  \|\varphi\|_\infty\,r^{-k}
  \int_{E\cap B(x,Rr)}|\theta(y)-\theta(x)|\,d\mathcal H^k(y)
  \longrightarrow0.
 \end{split}
\]
由无权放大极限便得到第一式。

一般的 \(E\) 未必局部面积有限。
如\cref{b1:prop:rectifiable-weight-density}的证明，
把它分解为可数个互不相交、有限面积的 Borel 片 \(E_\ell\)，
再按位置和权重截断。
对 \(\mu\)-几乎每个 \(x\in E_\ell\)，测度微分定理给出
\[
 \frac{\mu(B(x,r)\setminus E_\ell)}{\mu(B(x,r))}\longrightarrow0,
\]
而同一处带权密度满足
\(\mu(B(x,r))/(\omega_kr^k)\longrightarrow\theta(x)\in(0,\infty)\)。
因而片外部分为 \(o(r^k)\)；
在包含测试函数支集的 \(B(x,Rr)\) 上使用同一估计，
即可把片上的极限推广到整个 \(\mu\)。

最后设 \(c_j(\eta_{x,r_j})_\#\mu\to\nu\neq0\)，并令
\(\lambda=\theta(x)\mathcal H^k\llcorner P_x\)。
已知 \(r_j^{-k}(\eta_{x,r_j})_\#\mu\to\lambda\)。
取非负测试函数 \(\varphi\)，使 \(\int\varphi\,d\lambda>0\)。
若 \(c_jr_j^k\) 无界，则对应子列在这个测试函数上的积分无界，与淡收敛矛盾。
若它有趋零子列，则每个测试函数的积分沿该子列趋零，与 \(\nu\neq0\) 矛盾。
故可再取子列，使 \(c_jr_j^k\to c\in(0,\infty)\)，于是 \(\nu=c\lambda\)。
反向包含由第一式乘以任意正常数得到。
\end{proof}

这是一条几乎处处的结论。
两条相交直线上的长度测度在交点放大后仍是两条直线的长度测度之和；
它不是单一平面测度，但交点对长度测度为零，不违反定理。
反过来，由正有限密度推出可整流性，需要%
\cref{b1:thm:preiss-rectifiability}中的深层结果。
本节只把已经可整流的对象的局部放大写清楚，不把这个反向过程当成紧性抽取的直接后果。

切测度只记录质量分布，没有记录曲面的取向。
即使知道某点的全部切测度，也不能从中恢复曲面应当朝哪个方向积分。
接下来引入的流把这种取向信息与边界运算放在同一个线性框架中。
